% How to build a network: the ideas, the motifs and four worked examples. Included by nn_builder_manual.tex after "Designing a network,
% step by step". Every example is a file in examples/ (build_*.nnb) that the tests build with these parameter counts
% (crates/core/tests/text.rs, manual_examples_build), and every short text below builds too (manual_snippets_build).
\section{How to build a network: the ideas, then four worked examples}
\label{sec:build}

The step-by-step section above builds one network click by click. This section is about \emph{deciding what to build}: the few ideas
every design follows, the handful of small groups of blocks (``motifs'') that almost every network is made of, and four complete
examples for different jobs, each built from those motifs with the reasoning written down. Each example is given as text (the NNB
language, Section~\ref{sec:text}): paste it into the \ui{Text} tab, or open the file from \file{doc/examples/} with \ui{File}
$\rightarrow$ \ui{Import Design}, and the canvas draws it. The shapes and parameter counts printed here are the application's own.

\subsection{Every network has four parts}
\label{sec:build-parts}
Whatever its name, a network for images is four parts in a row:

\begin{figure}[htbp]\centering
\begin{tikzpicture}[node distance=7mm]
  \node[blk, fill=nnin!15, text width=2.6cm, align=center] (in) {\textbf{Input}\\[1pt]the data's shape\\(set by the Dataset tab)};
  \node[blk, fill=nnconv!15, text width=3.2cm, align=center, right=of in] (body) {\textbf{Body}\\[1pt]finds patterns: small ones first,\\larger ones from them};
  \node[blk, fill=nnpool!15, text width=2.7cm, align=center, right=of body] (br) {\textbf{Bridge}\\[1pt]turns the body's result into\\what the head needs};
  \node[blk, fill=nnout!15, text width=2.6cm, align=center, right=of br] (head) {\textbf{Head}\\[1pt]gives the answer\\(Output block)};
  \foreach \a/\b in {in/body,body/br,br/head} \draw[arr] (\a) -- (\b);
  \node[below=1mm of body, font=\scriptsize\sffamily, text width=3.4cm, align=center] {Conv2D, BatchNorm, Pool, residual blocks, attention, \dots};
  \node[below=1mm of br, font=\scriptsize\sffamily, text width=2.9cm, align=center] {Global pool, Flatten, To tokens, Token pool, Upsample, \dots};
  \node[below=1mm of head, font=\scriptsize\sffamily, text width=2.8cm, align=center] {Output, Segmentation output, Detection output};
\end{tikzpicture}
\caption{The four parts of every design. Most of the thinking (and most of the parameters) goes into the body; the bridge and the head
follow from the task.}
\label{fig:build-parts}
\end{figure}

\begin{description}[nosep]
  \item[Input.] Fixed by the data: \shape{3, 32, 32} for CIFAR-10, \shape{6, 128, 128} for six Sentinel-2 bands in 128-pixel patches.
        Choose the dataset first, so that every shape shown while building is the real one.
  \item[Body.] Where the network learns. Its job is to turn pixels into \emph{features}: numbers that say ``there is an edge here'',
        ``a texture like water here'', ``something shaped like a car here''. Early layers see a few pixels; every convolution and every
        halving of the size lets the next layer see a larger part of the image (its \emph{receptive field}, Section~\ref{sec:math}).
  \item[Bridge.] The body's output has the wrong shape for the answer, so one block changes it: Global pool or Flatten from an image to
        one vector, To tokens from an image to a sequence, Token pool from a sequence to one vector, Upsample / ConvTranspose2D back to the
        image's size for a mask.
  \item[Head.] The answer, matched to the task of the Dataset tab: \ui{Output} for classes (also multi-label and numbers),
        \ui{Segmentation output} for a class per pixel, \ui{Detection output} for boxes. A design has exactly one head (several only for
        multi-task data).
\end{description}

\subsection{What travels along the arrows}
\label{sec:build-kinds}
Every arrow carries one of four kinds of data, and each block accepts some kinds and gives one kind. Knowing which kind you hold, and
which block turns it into the next kind, is most of what building takes (Figure~\ref{fig:build-kinds}).

\begin{figure}[htbp]\centering
\begin{tikzpicture}[node distance=17mm, kind/.style={blk, minimum width=2.5cm, minimum height=11mm, align=center, font=\small\sffamily},
    keep/.style={font=\scriptsize\sffamily, align=center, text=black!65, text width=2.9cm}, on/.style={font=\scriptsize\sffamily}]
  \node[kind, fill=nnconv!12] (ts) {time series\\$(C, T, H, W)$};
  \node[kind, fill=nnconv!25, right=of ts] (img) {image\\$(C, H, W)$};
  \node[kind, fill=nnatt!15, right=of img] (tok) {tokens / sequence\\$(N, D)$};
  \node[kind, fill=nndense!15, right=of tok] (vec) {vector\\$(n)$};
  % what keeps each kind, under it
  \node[keep, below=1.5mm of ts] (kts) {Conv3D};
  \node[keep, below=1.5mm of img] (kimg) {Conv2D, Pool, BatchNorm, Add, Concat; ConvTranspose2D and Upsample grow it};
  \node[keep, below=1.5mm of tok] (ktok) {Self-attention, Dense, LayerNorm, LSTM, GRU, Conv1D};
  \node[keep, below=1.5mm of vec] (kvec) {Dense, Dropout, BatchNorm};
  % what turns one kind into another
  \draw[arr] (ts) -- node[above, on] {Temporal pool} (img);
  \draw[arr] (img) -- node[above, on] {To tokens} (tok);
  \draw[arr] (tok) -- node[above, on] {Token pool} (vec);
  \draw[arr] (img.north) to[out=40, in=140] node[above, on] {Global pool, Flatten} (vec.north);
  \draw[arr] (ts.north) .. controls +(0,22mm) and +(0,22mm) .. node[above, on] {To sequence (one vector per date)} (tok.north);
  % the heads
  \node[blk, fill=nnout!20, below=9mm of kvec, align=center] (out) {\textbf{Output}\\classes, labels, numbers};
  \node[blk, fill=nnout!20, below=9mm of kimg, align=center] (seg) {\textbf{Segmentation / Detection output}\\a class per pixel, boxes per cell};
  \draw[arr] (kvec) -- (out);
  \draw[arr] (kimg) -- node[right, on] {at the mask's size} (seg);
\end{tikzpicture}
\caption{The four kinds of data inside a network, the blocks that keep each kind (written under it), and the blocks that turn one kind into
another (on the arrows). There is no way back from a vector to an image or from tokens to an image: plan the order so you never need one.}
\label{fig:build-kinds}
\end{figure}

Two consequences explain most red blocks. A block must receive the kind it works on: Dense and Output need a vector (``add Flatten or
Global pool before it''), Conv2D needs an image, Self-attention needs tokens, and Output after To tokens needs a Token pool first. And
the head must match the task: a segmentation dataset wants an image at the mask's size at the end, not a vector.

\subsection{Seven questions to answer before the first block}
\label{sec:build-questions}
\begin{enumerate}[leftmargin=*]
  \item \textbf{What is the answer?} One class per image, several (multi-label), a few numbers, a class per pixel, or boxes. This fixes
        the head, and with it the bridge (Figure~\ref{fig:build-kinds}).
  \item \textbf{How large are the images, and how small must they get?} A classifier usually shrinks the image until it is about
        4--8 pixels across before Global pool, halving it each time: 32 pixels need 2--3 halvings, 64 pixels 3--4, 224 pixels 5. Each
        halving is a Pool or a convolution with stride 2. A segmentation network shrinks the same way and then grows back, one step for
        each halving.
  \item \textbf{How many channels?} Start with 16--64 and double the channels whenever the size halves, so every stage costs about
        the same: \shape{32, 32, 32} $\rightarrow$ \shape{64, 16, 16} $\rightarrow$ \shape{128, 8, 8}. The variable \texttt{let w = 32}
        in the text makes the width one number to change.
  \item \textbf{How much data is there?} The number of parameters should grow with the number of training images. As a rough guide
        from practice, not a law: a few thousand images suit up to a few hundred thousand parameters (or a large network started from
        ImageNet weights, Section~\ref{sec:train}); tens of thousands of images suit 0.3--5 million; more needs more. A network much
        larger than its data learns the training images by heart: training accuracy far above validation accuracy.
  \item \textbf{Convolutions or a Transformer?} Convolutions assume that nearby pixels belong together, which is true of photos and
        satellite images, and they learn well from little data. Self-attention relates every position to every other, which helps with
        long-range context, but it needs much more data to learn what convolutions know from the start. On small data, use convolutions,
        or convolutions first and a few attention blocks after them (Example~2).
  \item \textbf{How deep?} More than about ten convolutions in a row train badly without help. The help is BatchNorm (or GroupNorm for
        small batches) after each convolution, and residual connections: an \ui{Add} that adds a block's input to its output, so each
        block only has to learn a correction.
  \item \textbf{How will it be checked?} Plan a validation part (Train tab, 10--20\,\%) and early stopping from the first run. Compare a
        small design with a larger one before making it larger still: if training accuracy is low, the network is too small or too
        shallow (TEST~25 in Section~\ref{sec:tests}: 25\,\% on CIFAR-10 with 1\,898 parameters); if training accuracy is far above
        validation accuracy, it is too large for the data or needs more augmentation and dropout.
\end{enumerate}

\subsection{The motifs: small groups that networks are made of}
\label{sec:build-motifs}
Famous networks look complicated, but they are a few small groups of blocks repeated with different numbers. Learn these and almost
every design in this manual reads as a short sentence. Each motif below is a function (\texttt{def}, Section~\ref{sec:text}): written
once, then used like a block, with its numbers as arguments.

\paragraph{1. The convolution unit} (image $\rightarrow$ image). Convolution, BatchNorm, ReLU. The convolution has no bias and no activation
of its own: BatchNorm's shift replaces the bias, and the activation comes after the normalisation. With \texttt{s = 2} it halves the size
(a $4\times4$ kernel with stride 2 and padding 1 halves it exactly).
\begin{verbatim}
def unit(c, s = 1) {
  a conv c kernel=s+2 stride s pad 1 none nobias
  b bn
  r relu
}
2 unit(32)            # (3, 32, 32) -> (32, 32, 32)
3 unit(64, 2)         # -> (64, 16, 16)
\end{verbatim}

\paragraph{2. The stage} (image $\rightarrow$ smaller image). One halving and a few units at the new size. A classifier's body is two to
five stages, each with twice the channels of the one before. The simplest stage is a Pool and two units:
\begin{verbatim}
def unit(c) {
  a conv c 3x3 none nobias
  b bn
  r relu
}
def stage(c) {
  p maxpool 2x2
  x unit(c)
  y unit(c)
}
2 unit(32)
3 for c in 64, 128 { s stage(c) }     # (32, 32, 32) -> (64, 16, 16) -> (128, 8, 8)
\end{verbatim}

\paragraph{3. The residual block} (image $\rightarrow$ image). Two convolutions whose result is \emph{added} to the block's input, so the
block learns a correction to what it receives. When the block changes the size or the channels, the input is first brought to the new
shape by a \emph{shortcut} (here an average pool and a $1\times1$ convolution). On the canvas, \ui{Residual block after selected} inserts
one; the skip-connection tool adds a shortcut to blocks you already have.
\begin{figure}[htbp]\centering
\begin{tikzpicture}[node distance=4mm]
  \node[blk, fill=nnin!12] (x) {input};
  \node[blk, fill=nnconv!15, right=of x] (c1) {Conv $s$};
  \node[blk, fill=nnnorm!15, right=of c1] (b1) {BN};
  \node[blk, fill=nnpool!12, right=of b1] (r1) {ReLU};
  \node[blk, fill=nnconv!15, right=of r1] (c2) {Conv 3$\times$3};
  \node[blk, fill=nnnorm!15, right=of c2] (b2) {BN};
  \node[blk, fill=nnjoin!15, right=8mm of b2] (add) {Add};
  \node[blk, fill=nnpool!12, right=of add] (r2) {ReLU};
  \node[blk, fill=nnpool!12, below=6mm of c1] (p) {AvgPool $s$};
  \node[blk, fill=nnconv!15, right=of p] (k) {Conv 1$\times$1};
  \node[blk, fill=nnnorm!15, right=of k] (kn) {BN};
  \foreach \a/\b in {x/c1,c1/b1,b1/r1,r1/c2,c2/b2,b2/add,add/r2,p/k,k/kn} \draw[arr] (\a) -- (\b);
  \draw[arr] (x.south) |- (p.west);
  \draw[arr] (kn.east) -| (add.south);
  \node[font=\scriptsize\sffamily, below=0mm of kn.south east, anchor=north west, text=black!60] {shortcut: only needed when size or channels change};
\end{tikzpicture}
\caption{The residual block of Example~1 (without its gate). With $s=2$ the main path halves the size, and the shortcut does too.}
\label{fig:build-res}
\end{figure}

\paragraph{4. The channel gate} (image $\rightarrow$ image, same shape). Squeeze-and-excitation~\cite{hu2017senet}: Global pool
summarises each channel in one number, two small Dense layers decide how much each channel matters for this image (a sigmoid gives 0 to 1),
and Multiply scales the channels by it. Cheap (a few thousand parameters) and often worth one to two points of accuracy.
\begin{verbatim}
def gate(c) {
  g global_pool
  d dense c // 4
  e dense c sigmoid
  m mul <- in, e          # (c, H, W) times (c): each channel scaled
}
2 conv 64 3x3
3 gate(64)
\end{verbatim}

\paragraph{5. The Transformer encoder block} (tokens $\rightarrow$ tokens). LayerNorm, Self-attention, Add; LayerNorm, a two-layer MLP,
Add (the figure in ``Transformers in NN Builder'', Section~\ref{sec:transformers}). Attention lets every token look at every other; the MLP works on each token alone. Before the
first block, To tokens turns the image into tokens and Class token + position adds the class token and the positions; after the last,
a LayerNorm and Token pool give the vector. The width $D$ must divide by the number of heads.
\begin{verbatim}
let d = 64
2 conv d 4x4 stride 4 valid none   # patches
3 to_tokens patches
4 cls_pos
5 repeat 2 {
    a ln
    b attention heads 4
    c add <- in, b
    e ln
    f dense d * 2 gelu
    g dense d none
    h add <- c, g
  }
6 ln
7 token_pool cls
8 output
\end{verbatim}

\paragraph{6. The decoder step} (smaller image + skip $\rightarrow$ larger image). For segmentation: double the size (Upsample or
ConvTranspose2D), join the encoder's map of that size with Concat (the \emph{skip connection}, which brings back the fine detail that
shrinking lost), then units. U-Net~\cite{ronneberger2015unet} is an encoder of stages and a decoder of these steps. A function with two
inputs calls the second one \texttt{in2}:
\begin{verbatim}
def up(c) {
  u upsample x2 bilinear
  j concat <- u, in2
  k conv c 3x3
}
2 conv 16 3x3
3 maxpool 2x2
4 conv 32 3x3
5 up(16) <- 4, 2       # (32, H/2, W/2) and (16, H, W) -> (16, H, W)
6 seg_output
\end{verbatim}

\paragraph{7. Joining two inputs} (two images $\rightarrow$ one). Two band groups (Sentinel-2 and Sentinel-1, or the same place at two
dates) start from two Input blocks. \emph{Different} data gets its own branch, joined by Concat (fusion, Section~\ref{sec:fusion}).
\emph{The same} data at two dates should be described in the same terms: the second branch reuses the first branch's layers with
\ui{Shared weights}, and \ui{Subtract} (absolute) compares them (Example~4).

\subsection{From the questions to a design}
\label{sec:build-table}
\begin{longtable}{@{}p{5.2cm}p{9.8cm}@{}}
\toprule
\textbf{If you want \dots} & \textbf{build \dots}\\
\midrule\endhead
one class per photo, little data & 3--4 stages of units, Global pool, Dropout, Output; or a demo started from ImageNet weights\\
one class per photo, more accuracy & the same with residual blocks (Example~1); add channel gates; widen with \texttt{let w}\\
to use attention on images & a convolution stem down to 8--16 positions across, then To tokens and 2--6 encoder blocks (Example~2)\\
a class for every pixel & an encoder of stages and a decoder of steps with skips, ending at the mask's size (Example~3)\\
to compare two dates & a twin branch with shared weights and Subtract (Example~4)\\
two sensors together & one Input per band group, a branch each, Concat (Section~\ref{sec:fusion})\\
several labels per image, or numbers & the same body as for one class; only the Dataset tab's task changes (the Output adapts)\\
images over time & Conv3D or Temporal pool first (Section~\ref{sec:timeseries}), then any image body\\
a layer that is not in the list & a Custom block in PyTorch (Section~\ref{sec:custom})\\
\bottomrule
\end{longtable}

%-------------------------------------------------------------------------------------------------------------------------------------
\subsection{Example 1: a classifier for satellite photos}
\label{sec:build-ex1}
\paragraph{The job.} EuroSAT RGB~\cite{helber2019eurosat} (Hugging Face \texttt{timm/eurosat-rgb}): 27\,000 Sentinel-2 photos of
$64\times64$ pixels in 10 land-cover classes (forest, river, highway, \dots). One class per photo.

\paragraph{The decisions.}
\begin{itemize}[nosep]
  \item Head and bridge: one class per image $\rightarrow$ Output, after Global pool.
  \item Size: 64 pixels to 8 needs three halvings. The first is done at once by the stem (a unit with stride 2): the finest details of a
        10-metre pixel matter less than the texture of a field or a town, and every later layer costs a quarter as much.
  \item Body: three stages of one residual block each (32, 64, 128 channels at $32\times32$, $16\times16$, $8\times8$), each with a
        channel gate, since land cover is very much about \emph{which} textures are present.
  \item Size of the whole: 27\,000 images suit a few hundred thousand parameters: about 0.4 million here. Dropout 0.2 before the Output.
\end{itemize}

\begin{figure}[htbp]\centering
\begin{tikzpicture}[node distance=3.2mm, every node/.style={font=\scriptsize\sffamily, align=center}]
  \node[blk, fill=nnin!15] (a) {Input\\[-1pt]{\tiny\color{black!60}(3,64,64)}};
  \node[blk, fill=nnconv!15, right=of a] (b) {unit 32, $s$2\\[-1pt]{\tiny\color{black!60}(32,32,32)}};
  \node[blk, fill=nnjoin!12, right=of b] (c) {res 32\\[-1pt]{\tiny\color{black!60}(32,32,32)}};
  \node[blk, fill=nnjoin!12, right=of c] (d) {res 64, $s$2\\[-1pt]{\tiny\color{black!60}(64,16,16)}};
  \node[blk, fill=nnjoin!12, right=of d] (e) {res 128, $s$2\\[-1pt]{\tiny\color{black!60}(128,8,8)}};
  \node[blk, fill=nnpool!15, right=of e] (f) {Global pool\\[-1pt]{\tiny\color{black!60}(128)}};
  \node[blk, fill=nnpool!8, right=of f] (g) {Dropout\\[-1pt]{\tiny\color{black!60}(128)}};
  \node[blk, fill=nnout!15, right=of g] (h) {Output\\[-1pt]{\tiny\color{black!60}(10)}};
  \foreach \x/\y in {a/b,b/c,c/d,d/e,e/f,f/g,g/h} \draw[arr] (\x) -- (\y);
\end{tikzpicture}
\caption{Example 1 at the level of motifs (``res'' is the residual block of Figure~\ref{fig:build-res} with a channel gate); the
shapes are what each part gives. On the canvas it is 49 blocks.}
\end{figure}

\VerbatimInput[fontsize=\small,frame=single,numbers=left,numbersep=6pt,label=\file{build\_classifier.nnb}]{examples/build_classifier.nnb}

\paragraph{Reading it.} The three functions are the motifs 1, 4 and 3. The line labelled 3 is the body: a \texttt{for} over (channels, stride) that
calls \texttt{res} three times, each pass taking the previous pass's output. A residual block with stride 1 and unchanged channels would
not need its shortcut convolution; it is kept so that one function serves every stage (it costs 1\,088 parameters).

\paragraph{The numbers.} 393\,122 parameters, 58 M multiply--accumulates per image. Where they are:
\begin{center}
\begin{tabular}{@{}lrr@{}}
\toprule
\textbf{Part} & \textbf{Output} & \textbf{Parameters}\\
\midrule
stem: unit 32, stride 2 & \shape{32, 32, 32} & 1\,600\\
stage 1: residual block 32 + gate & \shape{32, 32, 32} & 20\,200\\
stage 2: residual block 64, stride 2 + gate & \shape{64, 16, 16} & 74\,192\\
stage 3: residual block 128, stride 2 + gate & \shape{128, 8, 8} & 295\,840\\
Global pool, Dropout, Output & \shape{10} & 1\,290\\
\midrule
total & & 393\,122\\
\bottomrule
\end{tabular}
\end{center}
Three quarters of the parameters sit in the last stage, as usual: its convolutions have the most channels. The first stage has few
parameters but the most computation (it works on $32\times32$ positions).

\paragraph{Training it.} \ui{Use recommended} (AdamW, cosine schedule with warm-up, label smoothing), 30--50 epochs, batch 128,
validation 10\,\% with early stopping; augmentation: horizontal \emph{and} vertical flips (a satellite photo has no ``up''). Kaggle, one
GPU: it is small.

\paragraph{Changing it.} Wider: \texttt{let w = 48}. Deeper: a second call in the loop, \texttt{y res(c)} after \texttt{x res(c, s)}
(this is the Small ResNet of Section~\ref{sec:text}). More bands (EuroSAT-MS, 13 bands): nothing changes in the text; only the Input's
shape and the stem's parameters grow. Larger images (e.g.\ 128 pixels): one more stage, \texttt{(w * 8, 2)}.

%-------------------------------------------------------------------------------------------------------------------------------------
\subsection{Example 2: convolutions, then a Transformer}
\label{sec:build-ex2}
\paragraph{The job.} CIFAR-10: 50\,000 photos of $32\times32$ pixels, 10 classes. This example grows TEST~25 (Section~\ref{sec:tests}:
one patch convolution, To tokens, Token pool, 25\,\% accuracy) into a network that uses its tokens.

\paragraph{The decisions.}
\begin{itemize}[nosep]
  \item Why TEST~25 stayed at 25\,\%: To tokens only reshapes, and Token pool (mean) averages the 64 patch tokens straight away. Nothing
        ever relates one patch to another, and a single $4\times4$ convolution sees only its own patch. The network is too \emph{simple},
        not too small: its training accuracy was as low as its test accuracy.
  \item The fix is two things. Convolutions first (a stem), so that each position already describes its neighbourhood: two units at
        $32\times32$, a Pool, one more unit, then a $2\times2$ convolution with stride 2 that makes an $8\times8$ grid of 128-value
        positions. Then encoder blocks, so that positions exchange information: four, width 128, 4 heads.
  \item Bridge: To tokens (64 tokens), Class token + position (65 tokens: attention ignores order unless positions are added); at the
        end LayerNorm and Token pool (class token).
  \item Size: 0.72 million parameters for 50\,000 images.
\end{itemize}

\begin{figure}[htbp]\centering
\begin{tikzpicture}[node distance=3.2mm, every node/.style={font=\scriptsize\sffamily, align=center}]
  \node[blk, fill=nnin!15] (a) {Input\\[-1pt]{\tiny\color{black!60}(3,32,32)}};
  \node[blk, fill=nnconv!15, right=of a] (b) {unit 64 $\times$2\\[-1pt]{\tiny\color{black!60}(64,32,32)}};
  \node[blk, fill=nnpool!15, right=of b] (c) {Pool\\[-1pt]{\tiny\color{black!60}(64,16,16)}};
  \node[blk, fill=nnconv!15, right=of c] (d) {unit 128\\[-1pt]{\tiny\color{black!60}(128,16,16)}};
  \node[blk, fill=nnconv!25, right=of d] (e) {Conv 2$\times$2 $s$2\\[-1pt]{\tiny\color{black!60}(128,8,8)}};
  \node[blk, fill=nntok!15, right=of e] (f) {To tokens\\[-1pt]{\tiny\color{black!60}(64,128)}};
  \node[blk, fill=nntok!15, right=of f] (g) {cls + pos\\[-1pt]{\tiny\color{black!60}(65,128)}};
  \node[blk, fill=nnatt!15, right=of g] (h) {encoder $\times$4\\[-1pt]{\tiny\color{black!60}(65,128)}};
  \node[blk, fill=nntok!15, right=of h] (i) {LN, Token pool\\[-1pt]{\tiny\color{black!60}(128)}};
  \node[blk, fill=nnout!15, right=of i] (j) {Output\\[-1pt]{\tiny\color{black!60}(10)}};
  \foreach \x/\y in {a/b,b/c,c/d,d/e,e/f,f/g,g/h,h/i,i/j} \draw[arr] (\x) -- (\y);
\end{tikzpicture}
\caption{Example 2: image kind on the left, tokens in the middle, a vector at the end (Figure~\ref{fig:build-kinds}).}
\end{figure}

\VerbatimInput[fontsize=\small,frame=single,numbers=left,numbersep=6pt,label=\file{build\_hybrid.nnb}]{examples/build_hybrid.nnb}

\paragraph{The numbers.} 718\,410 parameters, 101 M multiply--accumulates per image: the stem 178\,496 (with the $2\times2$
convolution), Class token + position 8\,448, each encoder block 132\,480 (attention 66\,048, MLP 65\,920, two LayerNorms 512), the
end 1\,546. Each line's shape can be checked against Figure~\ref{fig:build-kinds}: the block labelled 6 is the last of image kind, 7 the
first of token kind, 11 the first vector. (The labels, not the listing's grey line numbers.)

\paragraph{Training it.} \ui{Use recommended} (Transformers need the warm-up and weight decay more than CNNs do), 50--100 epochs,
batch 128, horizontal flip and random crop, validation 10\,\%. On one T4 an epoch should take seconds, not minutes; one GPU is enough.

\paragraph{Changing it.} To see what the stem is worth, delete the lines labelled 2--5 and make 6 \texttt{conv d 4x4 stride 4 valid none}: a
pure small ViT, which on CIFAR-10 from scratch usually does clearly worse. More encoder blocks: \texttt{repeat 6}. Wider: \texttt{let d
= 192} with \texttt{heads 6} or \texttt{heads 3} (192 must divide by the heads). A mean of all tokens instead of the class token:
\texttt{token\_pool mean} (and \texttt{cls\_pos nocls}).

%-------------------------------------------------------------------------------------------------------------------------------------
\subsection{Example 3: a water map from Sentinel-2}
\label{sec:build-ex3}
\paragraph{The job.} Six Sentinel-2 bands (B2, B3, B4, B8, B11, B12) in $128\times128$ patches, each with a mask of the same size:
water or not water for every pixel. Dataset tab: \ui{My own files}, task \ui{Segmentation}, images and masks folders
(Section~\ref{sec:data-kinds}).

\paragraph{The decisions.}
\begin{itemize}[nosep]
  \item Head: a class per pixel $\rightarrow$ Segmentation output, which must receive an image at the mask's size, $128\times128$.
  \item Body: an encoder that halves three times ($128 \rightarrow 16$), so that one value at the narrowest level sees $68\times68$ input
        pixels (a river's banks, a lake's shape; the inspector shows this receptive field), and a decoder that doubles three times back to 128. A skip connection at every size
        brings back the sharp edges.
  \item Motifs: \texttt{block} (two units) at every level, \texttt{down} (Pool + block) in the encoder, \texttt{up} (Upsample +
        Concat with the skip + block) in the decoder: motifs 1, 2 and 6.
  \item Bilinear Upsample instead of ConvTranspose2D: no parameters and no checkerboard patterns; the block after it learns what
        is needed. Dropout 0.3 at the narrowest point, where the network holds the most channels.
\end{itemize}

\begin{figure}[htbp]\centering
\begin{tikzpicture}[node distance=3mm, every node/.style={font=\scriptsize\sffamily, align=center}]
  \node[blk, fill=nnin!15] (in) {Input};
  \node[blk, fill=nnconv!15, right=of in] (e1) {block 16\\[-1pt]{\tiny\color{black!60}(16,128,128)}};
  \node[blk, fill=nnconv!15, below right=5mm and 2mm of e1] (e2) {down 32\\[-1pt]{\tiny\color{black!60}(32,64,64)}};
  \node[blk, fill=nnconv!15, below right=5mm and 2mm of e2] (e3) {down 64\\[-1pt]{\tiny\color{black!60}(64,32,32)}};
  \node[blk, fill=nnconv!25, below right=5mm and 2mm of e3] (e4) {down 128, Dropout\\[-1pt]{\tiny\color{black!60}(128,16,16)}};
  \node[blk, fill=nnjoin!12, above right=5mm and 2mm of e4] (d3) {up 64\\[-1pt]{\tiny\color{black!60}(64,32,32)}};
  \node[blk, fill=nnjoin!12, above right=5mm and 2mm of d3] (d2) {up 32\\[-1pt]{\tiny\color{black!60}(32,64,64)}};
  \node[blk, fill=nnjoin!12, above right=5mm and 2mm of d2] (d1) {up 16\\[-1pt]{\tiny\color{black!60}(16,128,128)}};
  \node[blk, fill=nnout!15, right=of d1] (out) {Seg output\\[-1pt]{\tiny\color{black!60}(2,128,128)}};
  \draw[arr] (in) -- (e1); \draw[arr] (e1) |- (e2); \draw[arr] (e2) |- (e3); \draw[arr] (e3) |- (e4);
  \draw[arr] (e4) -| (d3); \draw[arr] (d3) -| (d2); \draw[arr] (d2) -| (d1); \draw[arr] (d1) -- (out);
  \draw[arr, dashed] (e3) -- node[above] {skip} (d3);
  \draw[arr, dashed] (e2) -- node[above] {skip} (d2);
  \draw[arr, dashed] (e1) -- node[above] {skip} (d1);
\end{tikzpicture}
\caption{Example 3, a U-Net: down the left, up the right, a skip (Concat) across at every size.}
\end{figure}

\VerbatimInput[fontsize=\small,frame=single,numbers=left,numbersep=6pt,label=\file{build\_segment.nnb}]{examples/build_segment.nnb}

\paragraph{Reading it.} \texttt{up(64) <- 6, 4} gives the call two inputs: the line above's output (\texttt{in}, the dropout's
\shape{128, 16, 16}) and that of the line labelled 4 (\texttt{in2}, the encoder's \shape{64, 32, 32}). After Upsample the two have the same height and width,
so Concat can join them: \shape{128 + 64, 32, 32} = \shape{192, 32, 32}, which the block brings down to 64 channels.

\paragraph{The numbers.} 488\,450 parameters, 675 M multiply--accumulates per patch (the full-size levels are the expensive ones). The
encoder holds 294\,432 parameters, the decoder 193\,984, the Segmentation output 34 (a $1\times1$ convolution from 16 channels to 2
classes).

\paragraph{Training it.} \ui{Use recommended}, batch 16--32, horizontal and vertical flips (applied to image and mask alike), validation
10\,\% and early stopping on mIoU. Water is often a small part of each patch: look at \ui{Unbalanced classes}
(Section~\ref{sec:imbalance}) and at the per-class IoU in Results, not only at pixel accuracy.

\paragraph{Changing it.} Smaller objects or larger patches: one more level (\texttt{down(256)} and \texttt{up(128)}). Fewer
parameters: start at 8 channels. More classes (a land-cover map): only the Dataset tab changes; Segmentation output follows the classes.
Sentinel-1 radar as well: a second Input and a branch, joined at the bottleneck with Concat (Section~\ref{sec:fusion}).

%-------------------------------------------------------------------------------------------------------------------------------------
\subsection{Example 4: what changed between two dates}
\label{sec:build-ex4}
\paragraph{The job.} Pairs of Sentinel-2 patches of the same place (4 bands, $64\times64$) before and after: has the land changed
(cleared, built on, flooded) or not? The Dataset tab gives the ``before'' folder as band group 1 and the ``after'' folder as band group 2
with the same bands (Section~\ref{sec:siamese}).

\paragraph{The decisions.}
\begin{itemize}[nosep]
  \item Two Inputs (one per band group). Both images should be described \emph{in the same terms}, otherwise their difference measures
        the two branches, not the land. So the second branch uses the first branch's layers: \ui{Shared weights}. Build the ``before''
        branch, select its last block, press \ui{Twin branch (shared weights)}, and the ``after'' branch appears.
  \item The comparison: \ui{Subtract} with \emph{absolute difference}: large where the two descriptions differ, near zero where nothing
        changed.
  \item After it, one more convolution to judge the differences in context, then Global pool and Output (two classes).
\end{itemize}

\screenshot{build_change.png}{Example 4 on the canvas (drawn by the application; the right half is shown under the left). The upper branch's convolutions (2, 4, 6) are used
again by the lower branch's Shared-weights blocks (8, 10, 12); Subtract (13) compares the two at $16\times16$.}{fig:build-change}

\VerbatimInput[fontsize=\small,frame=single,numbers=left,numbersep=6pt,label=\file{build\_change.nnb}]{examples/build_change.nnb}

\paragraph{The numbers.} 241\,378 parameters. The twin branch adds none: its three Shared-weights blocks show ``the weights of another
block'' and 0 parameters, and the convolutions 2, 4 and 6 (93\,536 together) are trained by both images. Convolution 14 (147\,584) and the
Output (258) are the rest.

\paragraph{Training it.} \ui{Use recommended}, flips and turns (applied to both images alike), validation 10\,\%. Changed pairs are usually
rare: see Section~\ref{sec:imbalance}.

\paragraph{Changing it.} A change \emph{map} instead of one answer per pair: masks in the Dataset tab (segmentation), and after
Subtract a decoder of \texttt{up} steps to the image's size and a Segmentation output, as in Example~3 (the skips can come from either
branch). Keep both descriptions instead of their difference: \texttt{concat <- 6, 12} instead of the Subtract.

\subsection{When a design does not learn}
\label{sec:build-diagnose}
\begin{longtable}{@{}p{5.2cm}p{9.8cm}@{}}
\toprule
\textbf{What the Results show} & \textbf{What to change}\\
\midrule\endhead
training and validation accuracy both low, and close & the network is too simple or too small: add stages or blocks, widen it,
  bring attention or a stem (Example~2); check the receptive field covers the objects\\
training accuracy high, validation much lower & too large for the data: more augmentation, Dropout, weight decay, fewer channels;
  or ImageNet weights\\
the loss does not fall at all & learning rate too high or too low (try 0.001 with AdamW), missing normalisation in a deep network,
  or the data: run \ui{Check data} (Section~\ref{sec:train})\\
the loss is NaN & lower the learning rate, switch on gradient clipping (part of \ui{Use recommended})\\
``CUDA out of memory'' & smaller batch, \ui{Save memory}, fewer channels at the full-size levels\\
segmentation edges blurred & skip connections missing or joined at the wrong size; fewer halvings\\
\bottomrule
\end{longtable}
